% ---------------------------------------------------------------- ASSUNTO 2: ELETRICIDADE
\assunto{corB}{Eletricidade: circuitos e potência}{O assunto campeão do ENEM é circuitos elétricos: 11 questões em 5 anos, todo ano. Com potência e energia, são 14.}{Eletricidade}\label{pg:ele}

\secao{Parte 1 · As três grandezas e a Lei de Ohm}

\begin{analogia}[Eletricidade é como água no cano]
A caixa‑d'água lá no alto empurra a água pelo cano. Quanto \textbf{mais alta} a caixa, mais força a água tem. Se o cano tem um
trecho \textbf{estreito}, passa menos água. No circuito é igual: a pilha empurra, o fio conduz e o resistor aperta a passagem.\par
\smallskip{\small Só uma diferença: a água sai pela torneira, mas a corrente não é gasta. Ela só existe se o caminho for \textbf{fechado} e é a mesma antes e depois do resistor.}
\end{analogia}

\begin{center}
\begin{tikzpicture}
  % ---- lado da água
  \draw[tinta,line width=1pt,fill=corA!25] (0,2.6) rectangle (1.6,3.8);
  \fill[corA!60] (0.03,2.63) rectangle (1.57,3.45);
  \node[font=\small\bfseries,white] at (0.8,3.05) {caixa};
  \draw[tinta,line width=1pt] (0.4,2.6) -- (0.4,0.9) -- (2.0,0.9) -- (2.3,1.05) -- (3.3,1.05) -- (3.6,0.9) -- (5.0,0.9);
  \draw[tinta,line width=1pt] (0.9,2.6) -- (0.9,1.4) -- (2.0,1.4) -- (2.3,1.25) -- (3.3,1.25) -- (3.6,1.4) -- (5.0,1.4);
  \fill[corA!45] (0.42,2.6) -- (0.42,0.92) -- (2.0,0.92) -- (2.3,1.07) -- (3.3,1.07) -- (3.6,0.92) -- (5.0,0.92) -- (5.0,1.38) -- (3.6,1.38) -- (3.3,1.23) -- (2.3,1.23) -- (2.0,1.38) -- (0.88,1.38) -- (0.88,2.6) -- cycle;
  \draw[seta,white] (3.85,1.15) -- (4.8,1.15);
  \draw[cota] (-0.3,1.15) -- (-0.3,3.45);
  \node[font=\small,align=right,anchor=east] at (-0.45,2.3) {\selo{1}\ altura\\da caixa};
  \node[font=\small,anchor=north] at (2.8,0.8) {\selo{2}\ trecho estreito};
  \node[font=\small,anchor=south west,align=left] at (3.5,1.5) {\selo{3}\ água por\\segundo};
  % ---- ligação
  \node[font=\normalsize\bfseries,text=tema] at (7.15,2.0) {é como};
  % ---- lado elétrico
  \begin{scope}[shift={(9.0,0.7)}]
    \draw[fio] (0,0) to[battery1] (0,2.6) -- (1.0,2.6) to[R] (3.2,2.6) -- (4.0,2.6) -- (4.0,0) -- (0,0);
    \node[font=\small\bfseries,text=alerta] at (0.3,1.75) {$+$};
    \node[font=\small\bfseries,text=tinta] at (0.3,0.9) {$-$};
    \node[font=\small,anchor=north] at (0.4,-0.15) {\selo{1}\ tensão $U$ (pilha)};
    \node[font=\small,anchor=south] at (2.1,3.0) {\selo{2}\ resistência $R$};
    \draw[seta,tema] (3.7,2.25) -- (3.7,0.9);
    \node[font=\small,anchor=west] at (4.15,1.55) {\selo{3}\ corrente $i$};
  \end{scope}
\end{tikzpicture}
\end{center}

{\small\renewcommand{\arraystretch}{1.3}
\begin{tabular}{@{}p{2.6cm}p{0.9cm}p{2.1cm}>{\raggedright\arraybackslash}p{10.6cm}@{}}
\toprule
\textbf{Grandeza} & \textbf{Letra} & \textbf{Unidade} & \textbf{Na comparação com a água}\\
\midrule
Tensão & $U$ & volt (V) & a altura da caixa: o ``empurrão''. Também chamada de voltagem ou diferença de potencial (d.d.p.). Pilha: 1,5~V. Tomada: 127~V ou 220~V.\\
Corrente & $i$ & ampère (A) & quanta água passa por segundo.\\
Resistência & $R$ & ohm ($\Omega$) & o trecho estreito do cano: dificulta a passagem.\\
\bottomrule
\end{tabular}}

\begin{tcbraster}[fileira=3]
\begin{formula}{1ª Lei de Ohm}
{\Large $U = R \cdot i$}\par\medskip
{\small Macete: \textbf{URI}.\\ Mais resistência, menos corrente (para a mesma tensão).}
\end{formula}
\begin{cartao}[halign=flush center]
\begin{tikzpicture}[scale=0.8]
  \draw[tema,line width=1.2pt,fill=white] (0,0) -- (3,0) -- (1.5,2.3) -- cycle;
  \draw[tema,line width=1.2pt] (0.75,1.15) -- (2.25,1.15);
  \draw[tema,line width=1.2pt] (1.5,0) -- (1.5,1.15);
  \node[font=\Large\bfseries] at (1.5,1.6) {$U$};
  \node[font=\Large\bfseries] at (0.85,0.5) {$R$};
  \node[font=\Large\bfseries] at (2.15,0.5) {$i$};
\end{tikzpicture}\par
{\small \textbf{Tampe o que você quer achar.}\\ Tampou $R$: sobra $U \div i$.\\ Tampou $i$: sobra $U \div R$.}
\end{cartao}
\begin{formula}{2ª Lei de Ohm}
{\Large $R = \rho\,\dfrac{L}{A}$}\par\smallskip
\begin{tikzpicture}
  \draw[ouro,line width=5pt,line cap=round] (0.9,0.9) -- (2.3,0.9);
  \node[font=\footnotesize,anchor=north] at (1.6,0.75) {curto e grosso: passa fácil};
  \draw[ouro,line width=1.4pt,line cap=round] (0,0) -- (3.2,0);
  \node[font=\footnotesize,anchor=north] at (1.6,-0.08) {comprido e fino: resiste mais};
\end{tikzpicture}\par\smallskip
{\small $L$ é o comprimento, $A$ é a grossura e $\rho$ (``rô'') depende do material.}
\end{formula}
\end{tcbraster}

\begin{enemreal}{2022, questão 128}
\textbf{A situação.} Um aplicativo de trânsito compara cada rua a um resistor. A energia gasta por carro faz o papel da \textbf{tensão} ($U$).
Os carros por segundo fazem o papel da \textbf{corrente} ($i$). Rua mais congestionada é rua de \textbf{resistência maior}.
Qual é a ordem, da rota menos congestionada para a mais congestionada?\par\smallskip
\renewcommand{\arraystretch}{1.15}
\begin{tabular}{@{}llll@{}}
\textbf{\color{tema}Rota 1:} & $U = 510$, $i = 4$ & $\rightarrow$ & $R = 510 \div 4 = 127{,}5$\\
\textbf{\color{tema}Rota 2:} & $U = 608$, $i = 4$ & $\rightarrow$ & $R = 608 \div 4 = 152$\\
\textbf{\color{tema}Rota 3:} & $U = 575$, $i = 3$ & $\rightarrow$ & $R = 575 \div 3 \approx 192$\\
\end{tabular}
\resp{rotas 1, 2, 3 (da menor para a maior resistência)}
\nota{O ENEM gosta de vestir a Lei de Ohm com outra roupa. A conta é a mesma: $R = U \div i$.}
\end{enemreal}

\secao{Parte 2 · Série ou paralelo?}

\begin{tcbraster}[fileira=2]
\begin{quadro}{Em SÉRIE: um depois do outro}
\centering
\begin{circuitikz}
  \draw[fio] (0,0) to[battery1] (0,1.9) -- (0.6,1.9) to[lamp] (2.3,1.9) to[lamp] (4.0,1.9) -- (4.6,1.9) -- (4.6,0) -- (0,0);
  \draw[seta,tema] (3.2,-0.3) -- (1.4,-0.3) node[midway,below,font=\scriptsize,text=tinta] {um caminho só: mesma corrente};
\end{circuitikz}\par\smallskip
\raggedright\small
\textbf{No dia a dia:} pisca‑pisca antigo. Queimou uma lâmpada, apagam todas.
\begin{itemize}
  \item a \textbf{corrente é a mesma} em todos;
  \item a \textbf{tensão se divide} entre eles;
  \item $R_{\text{total}} = R_1 + R_2$ (a resistência aumenta).
\end{itemize}
\end{quadro}
\begin{quadro}{Em PARALELO: lado a lado}
\centering
\begin{circuitikz}
  \draw[fio] (0,0) to[battery1] (0,1.9) -- (4.2,1.9) to[lamp] (4.2,0) -- (0,0);
  \draw[fio] (2.2,1.9) to[lamp,*-*] (2.2,0);
  \draw[seta,tema] (0.5,2.2) -- (1.6,2.2);
  \draw[seta,tema] (1.85,1.6) -- (1.85,1.0);
  \draw[seta,tema] (3.85,1.6) -- (3.85,1.0);
  \node[font=\scriptsize,anchor=west] at (1.7,2.25) {a corrente se divide};
\end{circuitikz}\par\smallskip
\raggedright\small
\textbf{No dia a dia:} as tomadas da sua casa. Cada aparelho liga e desliga sozinho.
\begin{itemize}
  \item a \textbf{tensão é a mesma} em todos;
  \item a \textbf{corrente se divide} (a total é a soma);
  \item a resistência total \textbf{diminui}. Dois iguais: $R_{\text{total}} = R \div 2$.
\end{itemize}
\end{quadro}
\end{tcbraster}

\begin{analogia}[Paralelo é abrir mais caixas no supermercado]
Cada caixa novo é mais um caminho, e a fila toda anda mais rápido. Por isso a resistência total diminui.
E por isso, quanto mais aparelhos na mesma extensão, \textbf{maior a corrente total}.
\end{analogia}

\begin{minipage}[c]{0.40\linewidth}\centering
\begin{tikzpicture}
  \newcommand{\pilha}[3]{% x, y, sentido (1 = + à direita, -1 = + à esquerda)
    \draw[tinta,line width=0.8pt,fill=ouro!25] (#1,#2) rectangle (#1+1.5,#2+0.6);
    \fill[tinta] ({#1+0.75+#3*0.75},{#2+0.2}) rectangle ({#1+0.75+#3*0.87},{#2+0.4});
    \node[font=\small\bfseries,text=alerta] at ({#1+0.75+#3*0.5},{#2+0.3}) {$+$};
    \node[font=\small\bfseries] at ({#1+0.75-#3*0.5},{#2+0.3}) {$-$};}
  \pilha{0}{1.4}{1}\pilha{1.75}{1.4}{1}
  \node[anchor=west,font=\small\bfseries] at (3.6,1.7) {$1{,}5 + 1{,}5 = 3$~V};
  \pilha{0}{0.2}{1}\pilha{1.75}{0.2}{-1}
  \node[anchor=west,font=\small\bfseries,text=alerta] at (3.6,0.5) {$1{,}5 - 1{,}5 = 0$~V};
\end{tikzpicture}
\end{minipage}\hfill
\begin{minipage}[c]{0.57\linewidth}
\begin{macete}[Pilhas e baterias]
Em fila e no mesmo sentido (o $+$ de uma no $-$ da outra), as tensões \textbf{somam}. Uma pilha \textbf{invertida} joga contra: \textbf{subtrai}.
Caiu duas vezes em 2022 (questões 96 e 119).
\end{macete}
\end{minipage}

\begin{enemreal}{2025, questão 130}
\textbf{A situação.} Um aparelho de resistência $R_C$ precisa de uma certa tensão, mas a fonte disponível entrega 20\% a mais.
Como ligar um resistor de proteção $R_P$, e de quanto?
\passo{1. Série ou paralelo?} Em paralelo o aparelho continuaria recebendo a tensão toda. Em \textbf{série} a tensão se divide. É série.
\passo{2. Quanto?} Invente números: o aparelho precisa de 100~V e a fonte entrega 120~V. O resistor tem de ficar com os 20~V que sobram.
A corrente é a mesma nos dois e $U = R \cdot i$. Quem fica com 5 vezes menos tensão tem 5 vezes menos resistência: $R_P = R_C \div 5$.
\resp{em série, com $R_P = 0{,}2 \cdot R_C$}
\nota{A mesma ideia apareceu em 2024 (questão 131): um resistor em série com cada LED para ``segurar'' a tensão que sobra da bateria.}
\end{enemreal}

\secao{Parte 3 · Potência, energia e a conta de luz}

\begin{analogia}[Potência é a pressa de gastar energia]
O chuveiro (5\,500~W) gasta energia 550 vezes mais depressa que uma lâmpada de LED (10~W).
A conta de luz não cobra a potência. Cobra a \textbf{energia}, que é a potência multiplicada pelo \textbf{tempo ligado}.
\end{analogia}

\begin{tcbraster}[fileira=3]
\begin{formula}{Potência elétrica}
{\Large $P = U \cdot i$}\par\medskip
{\small Macete: \textbf{PIU}.\\ watt (W) = volt $\times$ ampère}
\end{formula}
\begin{formula}{Energia consumida}
{\Large $E = P \cdot \Delta t$}\par\medskip
{\small kWh = kW $\times$ horas.\\ É o que a conta de luz cobra.}
\end{formula}
\begin{formula}{Corrente que o aparelho puxa}
{\Large $i = \dfrac{P}{U}$}\par\medskip
{\small Disjuntor e fusível cortam a energia quando a corrente total passa do limite (por exemplo, 20~A).}
\end{formula}
\end{tcbraster}

\begin{exemplo}{Exemplo do dia a dia: quanto custa o banho?}
Chuveiro de 5\,500~W, banho de 15~min por dia, 30 dias, tarifa de R\$~1,00 por kWh.\par\smallskip
\begin{tabular}{@{}l@{\ \ }l@{}}
\textbf{1.} Potência em kW: & $5\,500 \div 1\,000 = 5{,}5$~kW\\
\textbf{2.} Tempo em horas: & 15~min $= 0{,}25$~h por dia. No mês: $0{,}25 \times 30 = 7{,}5$~h\\
\textbf{3.} Energia: & $E = 5{,}5 \times 7{,}5 = 41{,}25$~kWh\\
\textbf{4.} Custo: & $41{,}25 \times 1{,}00 =$ \textbf{R\$~41,25 por mês}\\
\end{tabular}
\end{exemplo}

\begin{macete}[Chuveiro: inverno ou verão?]
Na tomada, a tensão é fixa. Quem tem \textbf{menos resistência} puxa \textbf{mais corrente} e esquenta mais ($P = U^2 \div R$).
Chuveiro na posição ``inverno'' = resistência menor = potência maior. Na lâmpada, mais potência é mais brilho.
\end{macete}

\begin{enemreal}{2025, questão 94}
\begin{minipage}[c]{0.56\linewidth}
\textbf{A situação.} Uma régua de tomadas em 110~V tem fusível de 20~A. Ficam sempre ligados o computador (250~W) e o ar‑condicionado
portátil (1\,100~W). A estudante tenta usar, um de cada vez e nesta ordem: impressora (660~W), cafeteira (900~W), luminária (5~W) e
secador (750~W). Quantos desses aparelhos ela consegue usar?
\passo{1. Limite da régua:} $P = U \cdot i = 110 \times 20 = 2\,200$~W.
\passo{2. Já ligados:} $250 + 1\,100 = 1\,350$~W. Sobram 850~W.
\passo{3. Impressora:} $660 < 850$. Funciona.
\passo{4. Cafeteira:} $900 > 850$. O fusível queima e a régua para de funcionar. Acabou.
\resp{apenas 1}
\end{minipage}\hfill
\begin{minipage}[c]{0.41\linewidth}\centering
\begin{tikzpicture}[x=0.00285cm,y=1cm]
  % orçamento de potência: limite 2200 W
  \node[anchor=south west,font=\scriptsize\bfseries] at (0,1.5) {com a impressora};
  \fill[tinta!45] (0,1.05) rectangle (1350,1.5); \node[font=\scriptsize\bfseries,white] at (675,1.275) {1\,350 já ligados};
  \fill[corC] (1350,1.05) rectangle (2010,1.5); \node[font=\scriptsize\bfseries,white] at (1680,1.275) {660};
  \node[anchor=south west,font=\scriptsize\bfseries] at (0,0.45) {com a cafeteira};
  \fill[tinta!45] (0,0) rectangle (1350,0.45); \node[font=\scriptsize\bfseries,white] at (675,0.225) {1\,350 já ligados};
  \fill[alerta] (1350,0) rectangle (2250,0.45); \node[font=\scriptsize\bfseries,white] at (1800,0.225) {900};
  \draw[alerta,line width=1.2pt,dashed] (2200,-0.3) -- (2200,2.05);
  \node[font=\scriptsize\bfseries,text=alerta,anchor=south east] at (2260,2.05) {limite: 2\,200 W};
  \node[font=\scriptsize,anchor=north west,text=corC!70!black] at (1350,1.05) {cabe};
  \node[font=\scriptsize,anchor=north east,text=alerta] at (2260,-0.3) {passa do limite};
\end{tikzpicture}
\end{minipage}
\erro{continuar testando a luminária e o secador depois que o fusível já queimou.}
\end{enemreal}

\begin{enemreal}{2021, questão 128}
\textbf{A situação.} Um carro elétrico roda 5~km com cada kWh. Depois de uma viagem de 110~km, é recarregado num carregador de
220~V e 20~A. Quantas horas leva a recarga?
\passo{1. Energia gasta:} $110 \div 5 = 22$~kWh.
\passo{2. Potência do carregador:} $P = 220 \times 20 = 4\,400$~W $= 4{,}4$~kW.
\passo{3. Tempo:} $\Delta t = E \div P = 22 \div 4{,}4 = 5$~h.
\resp{5 horas}
\end{enemreal}

\begin{pegadinha}[As pegadinhas de eletricidade que o ENEM repete]
\begin{itemize}
  \item \textbf{kWh é energia, não potência.} E só dá certo com potência em \textbf{kW} e tempo em \textbf{horas}.
  \item \textbf{Aparelhos na mesma extensão estão em paralelo.} As potências (e as correntes) \textbf{somam}.
  \item \textbf{Mais resistência no caminho, menos corrente.} As lâmpadas que estão em série nesse caminho brilham \textbf{menos} (2023, questão 120).
  \item \textbf{Olhe a polaridade.} Pilha invertida subtrai. LED ligado ao contrário não acende.
  \item \textbf{Vai ordenar resistências?} Calcule $U \div i$ (ou $L \div A$) de cada uma antes de comparar. Não decida ``no olho''.
\end{itemize}
\end{pegadinha}
